\section{敏感性分析}

\subsection{为什么要把敏感性分析单独列出}

在上一节中，EVT 尾部扩展已经回答了“年度极端损失是否稳定、TVaR 99\% 是否需要单独校正”这一问题。但这还没有回答另一个同样重要的管理问题：\textbf{如果某类风险驱动上升，整体风险分布与尾部风险分别会怎样变化}。

因此，把敏感性分析单独列出会更清楚。它让报告结构分成两步：
\begin{enumerate}[leftmargin=2.2em]
  \item 先说明尾部该怎么建模，以及 EVT 结果是否稳健；
  \item 再说明不同风险驱动如何系统性改变均值、95\% 风险、99\% 风险与 EVT 尾部风险。
\end{enumerate}

这样一来，EVT 是方法增强，敏感性分析则是管理解释增强，二者职责清楚，不会在同一节里互相抢重点。

\subsection{敏感性分析设计}

为了进一步回答“哪些驱动最能放大整体风险与尾部风险”，扩展版在 \texttt{base} 情景之外增加了系统化的频率敏感性分析。做法很直接：保持严重度与 Markov 结构不变，只对特定风险块的 \textbf{频率乘子} 上调。

若某一风险块的基准频率均值为 $\lambda_b$，敏感性乘子为 $m_b$，则
\begin{equation}
\lambda_b^{\star} = m_b \lambda_b.
\end{equation}

本次敏感性分析包括两类设计：
\begin{itemize}[leftmargin=2.2em]
  \item \textbf{单因素冲击}：分别对 internal、system、external、cancel、divert 做 +10\%、+20\%、+30\% 上调；
  \item \textbf{组合压力}：对 \texttt{internal+system}、\texttt{external+system}、\texttt{internal+cancel}、\texttt{external+cancel+divert} 统一施加 +20\% 上调。
\end{itemize}

每个敏感性场景都重新进行 50{,}000 次年度模拟，并输出一组完整风险指标：
\begin{itemize}[leftmargin=2.2em]
  \item Expected impact；
  \item VaR 95\%、TVaR 95\%；
  \item VaR 99\%、TVaR 99\%；
  \item EVT VaR 99\%、EVT TVaR 99\%。
\end{itemize}

因此，这里并不是只对尾部做一张附加表，而是把 \textbf{均值、95\% 风险、99\% 风险与 EVT 尾部风险} 都放到同一个敏感性框架里观察。

\subsection{敏感性分析结果：从整体风险到尾部风险}

\begin{table}[H]
\centering
\caption{对 EVT TVaR 99\% 影响最大的敏感性场景}
\resizebox{\textwidth}{!}{
\begin{tabular}{lclrrr}
\toprule
场景 & 类型 & 驱动项 & Expected impact 增幅 & TVaR 99\% 增幅 & EVT TVaR 99\% 增幅 \\
\midrule
internal\_plus\_30pct & 单因素 & internal & 15.19\% & 17.08\% & 17.10\% \\
internal\_cancel\_plus\_20pct & 组合 & internal, cancel & 14.91\% & 15.57\% & 15.65\% \\
internal\_system\_plus\_20pct & 组合 & internal, system & 13.77\% & 14.46\% & 14.38\% \\
internal\_plus\_20pct & 单因素 & internal & 10.15\% & 10.99\% & 11.01\% \\
cancel\_plus\_30pct & 单因素 & cancel & 7.01\% & 7.29\% & 7.27\% \\
external\_cancel\_divert\_plus\_20pct & 组合 & external, cancel, divert & 6.18\% & 6.52\% & 6.73\% \\
internal\_plus\_10pct & 单因素 & internal & 5.02\% & 5.60\% & 5.62\% \\
cancel\_plus\_20pct & 单因素 & cancel & 4.85\% & 5.21\% & 5.29\% \\
\bottomrule
\end{tabular}}
\end{table}

\begin{figure}[H]
  \centering
  \includegraphics[width=0.9\textwidth]{chart_13_sensitivity_tornado.png}
  \caption{敏感性场景对 EVT TVaR 99\% 的影响排序}
  \label{fig:sensitivity-tornado}
\end{figure}

从表格和图 \ref{fig:sensitivity-tornado} 可以先提炼出最直接的尾部结论：
\begin{enumerate}[leftmargin=2.2em]
  \item \textbf{internal disruption 是最强的尾部放大器}。无论看单因素还是组合场景，internal 都稳定地排在最前列；
  \item \textbf{cancel 是第二层重要驱动}。取消频率的提升对 TVaR 99\% 有持续而明显的推升作用；
  \item \textbf{system 风险有影响，但不如 internal 直接}。它在组合压力中会显著放大尾部，但单独的边际影响弱于 internal；
  \item \textbf{external/divert 更像尾部辅助驱动}。它们能够抬高尾部风险，但在当前数据结构下并不是第一主导项。
\end{enumerate}

\begin{figure}[H]
  \centering
  \includegraphics[width=0.92\textwidth]{chart_14_sensitivity_heatmap.png}
  \caption{全指标敏感性矩阵：各场景相对基准的百分比变化}
  \label{fig:sensitivity-heatmap}
\end{figure}

图 \ref{fig:sensitivity-heatmap} 把所有核心风险指标放在同一张热力图里，因此可以同时回答两个问题：
\begin{enumerate}[leftmargin=2.2em]
  \item 某个驱动是只放大尾部，还是连均值和 95\% 风险也一起推高；
  \item 某个组合场景的影响是广泛性的，还是主要集中在 99\% 尾部。
\end{enumerate}

从这张图可以看出，\texttt{internal\_plus\_30pct}、\texttt{internal\_cancel\_plus\_20pct} 与 \texttt{internal\_system\_plus\_20pct} 不只是推高 EVT TVaR 99\%，也会同步抬高 Expected impact、VaR 95\%、TVaR 95\%、VaR 99\% 与 TVaR 99\%。这说明 internal 风险并不是一个“只在最极端年份才重要”的驱动，而是会系统性抬升整个年度损失分布。

\begin{figure}[H]
  \centering
  \includegraphics[width=0.92\textwidth]{chart_15_single_factor_sensitivity_profiles.png}
  \caption{单因素冲击下各核心指标的敏感性曲线}
  \label{fig:sensitivity-profiles}
\end{figure}

图 \ref{fig:sensitivity-profiles} 则专门看单因素冲击的轨迹。它表明：
\begin{enumerate}[leftmargin=2.2em]
  \item \textbf{internal} 的冲击曲线在几乎所有指标上都最陡，说明它对整体分布和尾部同时敏感；
  \item \textbf{cancel} 对尾部指标的影响强于对均值的影响，体现出取消更偏向 tail amplifier 的特征；
  \item \textbf{system} 有稳定但中等强度的影响；
  \item \textbf{external} 和 \textbf{divert} 的影响相对温和，在弱冲击下更像辅助驱动。
\end{enumerate}

\subsection{这一节对原项目意味着什么}

把敏感性分析单独列出之后，原项目的管理叙事也更清楚了：
\begin{enumerate}[leftmargin=2.2em]
  \item 它不只是在报告一组静态情景指标，而是在回答“\textbf{哪一类运营驱动最值得优先配置 safety buffer}”；
  \item 它不只是在解释 99\% 尾部，而是在同时解释“\textbf{均值、95\% 风险、99\% 风险和 EVT 尾部风险如何一起变化}”。
\end{enumerate}

因此，这一节的价值在于把模型结果从“风险测量”推进到“风险管理优先级排序”。这也是把敏感性分析单独列出来之后最明显的收益：结论更聚焦，读者更容易抓住你真正想强调的管理含义。
